%% modified version of the script by J. Leon, Beerware
\tikzsetnextfilename{fig_lstm_architecture}
\centering
\begin{tikzpicture}[
	% GLOBAL CFG
	font= \scriptsize,
	>=LaTeX,
	% Styles
	cell/.style={% For the main box
		rectangle, 
		rounded corners=5mm, 
		draw,
		very thick,
		fill=Prune,
		opacity=.75,
	},
	operator/.style={%For operators like +  and  x
		circle,
		draw,
		inner sep=-0.5pt,
		minimum height =.4cm,
		fill=white,
		opacity=.5,
		text opacity=1,
		thick,
	},
	function/.style={%For functions
		ellipse,
		draw,
		inner sep=1pt,
		fill=white,
		opacity=.5,
		text opacity=1,
		thick,
	},
	ct/.style={% For external inputs and outputs
		circle,
		draw,
		line width = .75pt,
		minimum width=1cm,
		inner sep=1pt,
	},
	gt/.style={% For internal inputs
		rectangle,
		draw,
		minimum width=4mm,
		minimum height=3mm,
		inner sep=1pt,
		fill=blue!30,
		opacity=.5,
		text opacity=1,
	},
	mylabel/.style={% something new that I have learned
		font=\scriptsize%\sffamily
	},
	ArrowC1/.style={% Arrows with rounded corners
		rounded corners=.25cm,
		thick,
	},
	ArrowC2/.style={% Arrows with big rounded corners
		rounded corners=.5cm,
		thick,
	},
	]
	
	%Start drawing the thing...    
	% Draw the cell: 
	\node [cell, minimum height =4cm, minimum width=6cm] at (0,0){} ;
	
	% Draw inputs named ibox#
	\node [gt] (ibox1) at (-2,-0.75) {$\boldsymbol\sigma$};
	\node [gt] (ibox2) at (-1.5,-0.75) {$\boldsymbol\sigma$};
	\node [gt, minimum width=.75cm] (ibox3) at (-0.5,-0.75) {$\tanh$};
	\node [gt] (ibox4) at (0.5,-0.75) {$\boldsymbol\sigma$};
	
	% Draw operators named mux# , add# and func#
	\node [operator] (mux1) at (-2,1.5) {\large$\boldsymbol\times$};
	\node [operator] (add1) at (-0.5,1.5) {\large$\boldsymbol+$};
	\node [operator] (mux2) at (-0.5,0) {\large$\boldsymbol\times$};
	\node [operator] (mux3) at (1.5,0) {\large$\boldsymbol\times$};
	\node [function] (func1) at (1.5,0.75) {$\tanh$};
	
	% Draw External inputs named as basis c,h,x
	\node[ct, label={[mylabel]Cellule}] (c) at (-4,1.5) {$\vect{c}_{t-1}$};
	\node[ct, label={[mylabel]Caché}] (h) at (-4,-1.5) {$\vect{h}_{t-1}$};
	\node[ct, label={[mylabel]left:Entrée}] (x) at (-2.5,-3) {$\vect{e}_t$};
	
	% Draw External outputs? named as basis c2,h2
	\node[ct] (c2) at (4,1.5) {$\vect{c}_t$};
	\node[ct, label={[mylabel]Sortie}] (h2) at (4,-1.5) {$\vect{h}_t$};
	
	% Start connecting all.
	%Intersections and displacements are used. 
	% Drawing arrows    
	\draw [ArrowC1] (c) -- (mux1) -- (add1) -- (c2);
	
	% Inputs
	\draw [ArrowC2] (h) -| (ibox4);
	\draw [ArrowC1] (h -| ibox1) ++(-0.5,0) -| (ibox1); 
	\draw [ArrowC1] (h -| ibox2)++(-0.5,0) -| (ibox2);
	\draw [ArrowC1] (h -| ibox3)++(-0.5,0) -| (ibox3);
	\draw [ArrowC1] (x) -- (x |- h)-| (ibox3);
	
	\draw[fill] (-2.25,-1.5) circle (2pt);
	
	% Internal
	\draw [->, ArrowC2] (ibox1) -- (mux1) 
		node[left, midway] {$\vect{f}_t$};
	\draw [->, ArrowC2] (ibox2) |- (mux2)
		node[midway] {$\vect{i}_t$};
	\draw [->, ArrowC2] (ibox3) -- (mux2)
		node[right, midway] {$\tilde{\vect{c}}_t$};
	\draw [->, ArrowC2] (ibox4) |- (mux3)
		node[midway] {$\vect{o}_t$};
	\draw [->, ArrowC2] (mux2) -- (add1);
	\draw [->, ArrowC1] (add1 -| func1)++(-0.5,0) -| (func1);
	\draw [->, ArrowC2] (func1) -- (mux3);
	
	%Outputs
	\draw [-, ArrowC2] (mux3) |- (h2);
	%\draw (c2 -| x2) ++(0,-0.1) coordinate (i1);
	%\draw [-, ArrowC2] (h2 -| x2)++(-0.5,0) -| (i1);
	%\draw [-, ArrowC2] (i1)++(0,0.2) -- (x2);
	
\end{tikzpicture}
\vspace{1mm}
\captionof{figure}[Schéma d'une architecture LSTM classique]{Architecture LSTM classique, les couches neuronales sont représentées par les rectangles violets, et les opérateurs élémentaires (opération élément par élément, par exemple produit de Hadamard pour les multiplications \parencite{styan_hadamard_1973}, noté $\hadamard$) sont représentés par des ellipses ou cercles gris. Le point noir indique la concaténation de $\vect{h}_{t-1}$ et $\vect{e}_t$.}
\label{fig_lstm_architecture}